\documentclass{article}
\usepackage{graphicx} % Required for inserting images
\usepackage{amsmath}
\title{Worksheet \#1}
\author{Rosa Orellana}
\date{August 3, 2026}

\begin{document}

\maketitle
\begin{enumerate}
\item Find all permutations $\pi\in S_4$ for which $\pi(f) = f$, where $f(x_1, x_2, x_3,x_4) = x_1x_2^2x_4^4 + x_2^2x_3^4x_4+x_1^4x_2^2x_3$. 

\item Find a polynomial $f$ in $x_1, x_2, x_3$ which has $\pi(f) = \text{sgn}(\pi) f$ for all $\pi\in S_3$, and which has $x_1^7x_3^2$ as one of its terms. Recall that $\text{sgn}(\pi)$ is the {\bf sign}
of the permutation that is 1, if $\pi$ is the product of an even number of transpositions and $-1$ if $\pi$ is the product of an odd number of permutations. 

\item Compute the four monomial symmetric polynomials $m_{3,2}(x_1,x_2)$,\  $m_{3,2}(x_1,x_2,x_3)$, $m_{3,3,1,1}(x_1,x_2,x_3)$, and $m_{3,3,1,1}(x_1,x_2,x_3,x_4)$. 

\item Compute the product of the formal power series $f(x) = \sum_{i = 0} ^\infty x^i$ and $g(x) = \sum_{j=0}^\infty jx^j$. 

\item Compute the product of the formal power series $f= \sum_{i=1}^\infty x_ix_{i+1}$ and $g=\sum_{j=1}^\infty x_j$.  

\item Compute $m_{2,1}$. Is $m_2m_1 = m_{2,1}$? 

\item Using symmetry, write $m_3m_2m_1$ as a linear combination of $\{m_\lambda\ : \ \lambda \vdash 6\}$.

\item Without multiplying everything out, find the coefficient of $m_{2,1^{a+b}}$ in the product $m_{1^{a+1}}m_{1^{b+1}}$. 

\item Write the following as a sum of monomial symmetric functions: 
\begin{enumerate}
\item $\prod_{i=1}^\infty (1+x_i)$.
\item $\prod_{i=1}^\infty (1-x_i)$
\end{enumerate}

\item Use the geometric series to write a very simple infinite product using the variables $x_i$'s that is equal to the sum 
\[ \sum_\lambda m_\lambda,\, \]
where here the sum is over all partitions of all positive integers. 
\end{enumerate}

\section*{Español by Ignacio Rojas}
\begin{enumerate}
\item Encuentre todas las permutaciones $\pi\in S_4$ para las cuales $\pi(f) = f$, donde $f(x_1, x_2, x_3,x_4) = x_1x_2^2x_4^4 + x_2^2x_3^4x_4+x_1^4x_2^2x_3$.
\item Encuentre un polinomio $f$ en $x_1, x_2, x_3$ que satisfaga $\pi(f) = \text{sgn}(\pi) f$ para toda $\pi\in S_3$, y que tenga $x_1^7x_3^2$ como uno de sus términos. Recuerde que $\text{sgn}(\pi)$ es el {\bf signo} de la permutación, que es 1 si $\pi$ es el producto de un número par de transposiciones y $-1$ si $\pi$ es el producto de un número impar de permutaciones.
\item Calcule los cuatro polinomios simétricos monomiales $m_{3,2}(x_1,x_2)$, $m_{3,2}(x_1,x_2,x_3)$, $m_{3,3,1,1}(x_1,x_2,x_3)$ y $m_{3,3,1,1}(x_1,x_2,x_3,x_4)$.
\item Calcule el producto de las series de potencias formales $f(x) = \sum_{i = 0} ^\infty x^i$ y $g(x) = \sum_{j=0}^\infty jx^j$.
\item Calcule el producto de las series de potencias formales $f= \sum_{i=1}^\infty x_ix_{i+1}$ y $g=\sum_{j=1}^\infty x_j$.
\item Calcule $m_{2,1}$. ¿Vale que $m_2m_1 = m_{2,1}$?
\item Usando la simetría, escriba $m_3m_2m_1$ como una combinación lineal de $\{m_\lambda\ : \ \lambda \vdash 6\}$.
\item Sin multiplicar todo, encuentre el coeficiente de $m_{2,1^{a+b}}$ en el producto $m_{1^{a+1}}m_{1^{b+1}}$.
\item Reescriba lo siguiente como una suma de funciones simétricas monomiales:
\begin{enumerate}
\item $\prod_{i=1}^\infty (1+x_i)$.
\item $\prod_{i=1}^\infty (1-x_i)$
\end{enumerate}
\item Utilizando la serie geométrica, escriba un producto infinito muy sencillo usando las variables $x_i$ que sea igual a la suma
\[ \sum_\lambda m_\lambda,\, \]
donde aquí la suma es sobre todas las particiones de todos los enteros positivos.
\end{enumerate}



\end{document}
